\documentclass{article}
\usepackage{amsmath}
\usepackage[ruled,vlined,linesnumbered]{algorithm2e}

\begin{document}

\section*{\textit{Divide et impera}}
Prima di discutere l'algoritmo ci tocca parlare della tecnica che quest'ultimo
sfrutta: \textbf{divide et impera}.
Caratterizzata da tre passaggi:
\begin{itemize}
  \item \textbf{divide}: il problema iniziale \textbf{P} di dimensione $n$ 
  viene suddiviso in $k$ sotto-problemi di dimensione più piccola di $n$,
  indicato come $\frac{n}{b}$.
  \item risoluzione \textbf{ricorsiva} dei sotto problemi; se la ricorsione
  arriva ad un caso base allora risolviamo direttamente.
  \item \textbf{impera}: combiniamo le soluzioni di ciasciun sotto-problema
  per ottenere la soluzione del problema originale \textbf{P}.
\end{itemize}

Il costo computazionale di un qualsiasi algoritmo \textit{divide et impera}
prende questa forma:
\[ T(n) = k \cdot T(\frac{n}{b}) + f(n) \]

Dove $k$ è il numero di sotto-problemi generati, $\frac{n}{b}$ la loro dimensione
e $f(n)$ il costo per dividire e combinare i problemi; questa formula prende la
stessa struttura di quella studiata nel capitolo sugli \textit{algoritmi ricorsivi},
poichè il \textit{divide et impera} è un applicazione diretta del paradigma di 
ricorsione.

\section*{Quicksort non in-loco}
\vspace{0.3cm}

\SetKwProg{Fn}{function}{}{}
\SetKw{Return}{\textbf{return}}

\begin{algorithm}[H]
\renewcommand{\thealgocf}{}
\caption{Ordina l'array}
\Fn{QUICKSORT(array $A$)}{
    x $\leftarrow$ elemento di $A$\;
    \BlankLine
    array $A_1$ $\leftarrow$ $\{ \text{elementi di A minori di } x\}$ \;
    array $A_2$ $\leftarrow$ $\{ \text{elementi di A maggiori di } x\}$ \;
    array $A_3$ $\leftarrow$ $\{ \text{valori di A uguali a } x\}$ \;
    \BlankLine
    \If{$|A_1| > 1$}{$QUICKSORT(A_1)$}
    \If{$|A_2| > 1$}{$QUICKSORT(A_2)$}
    \BlankLine
    copiamo $A_1$ e $A_2$ in $A$
}
\end{algorithm}
\vspace{0.3cm}

L'algoritmo \textbf{\textit{QuickSort}} è un algoritmo di \textbf{ordinamento}
che sftutta la metodologia \textit{divie et impera}:

\begin{itemize}
  \item \textbf{divide}: l'array in tanti sotto-array in base al valore del 
  \textit{perno} (ovvero la variabile $x$)
  \item risolve ricorsivamente $A_1$ e $A_2$ 
  \item \textbf{impera}: concatena le due sottosequenze ordinate
\end{itemize}

Il passaggio del \textbf{divide} è più complesso di 
quello \textbf{impera}, al contrario del \textbf{MergeSort}. \\
L'algoritmo però non è molto efficiente per un grande proble di costo in 
\textbf{memoria}, infatti ad ogni chiamata vengono allocati nuovi sottoarray
$A_1, A_2$ e $A_3$. \\
Per questo motivo è nato il \textbf{QuickSort in-loco}.

\section*{Quicksort in-loco}

\vspace{0.3cm}

\SetKwProg{Fn}{function}{}{}
\SetKw{Return}{\textbf{return}}

\begin{algorithm}[H]
\renewcommand{\thealgocf}{}
\caption{Ordina l'array, versione in-loco}

\Fn{QUICKSORT(array A, int i, int f)}{

  \If{ $i \ge f$ } {\Return \;}

  \BlankLine

  m $\leftarrow$ PARTITION(A, i, f) \;
  QUICKSORT($A$, $i$, $m - 1$) \;
  QUICKSORT($A$, $m + 1$, $f$) \;

}

\BlankLine

\Fn{ PARTITION(array A, int i, int f) } {

  x $\leftarrow$ A[i]\;
  inf $\leftarrow i$\;
  sup $\leftarrow f + 1$\;

  \BlankLine

  %TODO%
}

\end{algorithm}
\vspace{0.3cm}

\end{document}